\documentclass[a4paper,11pt]{article}

\newcommand{\TPnumero}{Lab 6 -- Solution}
% =========================================================
% Tutorial / lab preamble — Time Series Analysis module
% article class + fancyhdr + tcolorbox + listings (English)
% Author : Aymen Ben Brik (aymen.benbrik@esprit.tn)
% =========================================================

\usepackage[utf8]{inputenc}
\usepackage[T1]{fontenc}
\usepackage[english]{babel}
\usepackage{lmodern}
\usepackage{amsmath, amssymb, amsthm, mathtools, bm}
\usepackage{graphicx}
\usepackage{booktabs}
\usepackage{array}
\usepackage{multirow}
\usepackage{colortbl}
\usepackage{xcolor}
\usepackage{tikz}
\usetikzlibrary{positioning, arrows.meta, calc, shapes, fit, backgrounds, matrix}
\usepackage{listings}
\usepackage[most]{tcolorbox}
\usepackage{geometry}
\geometry{a4paper, margin=2cm, headheight=15pt}
\usepackage{fancyhdr}
\usepackage[hidelinks]{hyperref}
\usepackage{enumitem}

% --- Colours ---
\definecolor{espritBlue}{RGB}{30, 60, 110}
\definecolor{espritAccent}{RGB}{200, 80, 30}
\definecolor{boxEx}{RGB}{30, 60, 110}
\definecolor{boxQ}{RGB}{180, 120, 0}
\definecolor{boxC}{RGB}{30, 100, 60}
\definecolor{boxAtt}{RGB}{200, 80, 30}

% --- Listings ---
\definecolor{codebg}{RGB}{248,248,248}
\definecolor{codeframe}{RGB}{220,220,220}
\definecolor{keyword}{RGB}{0,82,155}
\definecolor{comment}{RGB}{88,110,117}
\definecolor{string}{RGB}{133,153,0}

\lstdefinestyle{pythonstyle}{
  language=Python,
  basicstyle=\ttfamily\small,
  keywordstyle=\color{keyword}\bfseries,
  commentstyle=\color{comment}\itshape,
  stringstyle=\color{string},
  backgroundcolor=\color{codebg},
  frame=single,
  rulecolor=\color{codeframe},
  frameround=tttt,
  breaklines=true,
  showstringspaces=false,
  tabsize=2
}
\lstdefinestyle{Rstyle}{
  language=R,
  basicstyle=\ttfamily\small,
  keywordstyle=\color{keyword}\bfseries,
  commentstyle=\color{comment}\itshape,
  stringstyle=\color{string},
  backgroundcolor=\color{codebg},
  frame=single,
  rulecolor=\color{codeframe},
  frameround=tttt,
  breaklines=true,
  showstringspaces=false,
  tabsize=2
}
\lstset{style=pythonstyle}

% --- Common macros (configurable path) ---
\providecommand{\macrosPath}{../../preamble/macros.tex}
\input{\macrosPath}

% --- Exercise counter ---
\newcounter{excounter}
\renewcommand{\theexcounter}{\arabic{excounter}}

\newtcolorbox[use counter=excounter]{exercise}[1][]{
  enhanced, breakable,
  colback=boxEx!5, colframe=boxEx, fonttitle=\bfseries,
  title={Exercise~\theexcounter\ifstrempty{#1}{}{ -- #1}},
  arc=2pt, boxrule=0.6pt
}
\newtcolorbox{question}[1][]{
  enhanced, breakable,
  colback=boxQ!5, colframe=boxQ, fonttitle=\bfseries,
  title={Question\ifstrempty{#1}{}{ -- #1}},
  arc=2pt, boxrule=0.5pt
}
\newtcolorbox{solution}[1][]{
  enhanced, breakable,
  colback=boxC!5, colframe=boxC, fonttitle=\bfseries,
  title={Solution\ifstrempty{#1}{}{ -- #1}},
  arc=2pt, boxrule=0.5pt
}
\newtcolorbox{warning}[1][]{
  enhanced, breakable,
  colback=boxAtt!5, colframe=boxAtt, fonttitle=\bfseries,
  title={Warning\ifstrempty{#1}{}{ -- #1}},
  arc=2pt, boxrule=0.5pt
}

% --- Headers / footers ---
\pagestyle{fancy}
\fancyhf{}
\fancyhead[L]{\textsc{\moduleNom} -- \auteurNom}
\fancyhead[R]{\TPnumero}
\fancyfoot[C]{\thepage}
\fancyfoot[L]{\auteurInstitution}
\fancyfoot[R]{\auteurEmail}
\renewcommand{\headrulewidth}{0.4pt}
\renewcommand{\footrulewidth}{0.2pt}

% --- Placeholder ---
\providecommand{\TPnumero}{Lab}


\title{Lab 6: ARCH effects, GARCH and Value-at-Risk\\
\large \emph{Reference solution}}
\author{\auteurNom\\\auteurEmail\\\auteurInstitution}
\date{\anneeUniversitaire}

\begin{document}
\maketitle

\noindent\emph{Reference Python code in \texttt{correction.py}.
The synthetic dataset is regenerated by \texttt{generate\_returns.py}
(GARCH$(1, 1)$ with $(\alpha_0, \alpha_1, \beta_1) = (0.05, 0.10,
0.85)$, $T = 1500$, seed $42$).}

\setcounter{excounter}{0}

% =========================================================
\begin{exercise}[Volatility clustering]\end{exercise}

\begin{solution}
$T = 1500$, sample mean $= +0.0454$, sample sd $= 0.9515$, excess
kurtosis $= 0.44$. The time plot shows quiet stretches alternating
with bursts. The ACF of $r_t$ stays inside the Bartlett band, but
the ACF of $r_t^2$ has many large positive spikes -- the
\emph{volatility clustering} signature. Returns are
serially uncorrelated but \emph{not} independent.
\end{solution}

% =========================================================
\begin{exercise}[Engle's LM test]\end{exercise}

\begin{solution}
\begin{center}
\begin{tabular}{rrrrl}
\toprule
lag $q$ & LM stat & $\chi^2_{q, 0.95}$ & $p$-value & Decision \\
\midrule
$5$    & $158.57$ & $11.07$ & $\approx 0$ & \textbf{reject} -- ARCH \\
$10$   & $162.86$ & $18.31$ & $\approx 0$ & \textbf{reject} -- ARCH \\
$20$   & $172.66$ & $31.41$ & $\approx 0$ & \textbf{reject} -- ARCH \\
\bottomrule
\end{tabular}
\end{center}
$LM \gg \chi^2_{q, 0.95}$ at every lag: strong evidence of ARCH
effects, consistent with the heavy ACF spikes on $r_t^2$.
\end{solution}

% =========================================================
\begin{exercise}[ARCH$(q)$ fitting]\end{exercise}

\begin{solution}
\begin{center}
\begin{tabular}{rrrrr}
\toprule
$q$ & log-lik. & AIC      & $\sum_i \widehat\alpha_i$ & First $\widehat\alpha_i$'s \\
\midrule
$1$ & $-2033.84$ & $4073.68$ & $0.159$ & $0.159$ \\
$3$ & $-2020.11$ & $4050.23$ & $0.299$ & $0.105,\,0.093,\,0.101$ \\
$5$ & $-2007.27$ & $\bm{4028.54}$ & $0.448$ & $0.078,\,0.081,\,0.085,\,0.107,\,0.096$ \\
\bottomrule
\end{tabular}
\end{center}
$\sum \widehat\alpha_i < 1$ in every case (covariance-stationary).
ARCH$(5)$ improves over ARCH$(3)$ but the gain is modest: the
spread of the loadings shows that pure ARCH cannot capture the
\emph{long memory} of the conditional variance with few parameters.
\end{solution}

% =========================================================
\begin{exercise}[GARCH$(1, 1)$ fitting]\end{exercise}

\begin{solution}
MLE estimates with constant mean and Gaussian innovations:
\[
  \widehat\mu = +0.0488,
  \quad
  \widehat\alpha_0 = 0.0878,
  \quad
  \widehat\alpha_1 = 0.1069,
  \quad
  \widehat\beta_1  = 0.7960.
\]
True DGP: $(0.05, 0.10, 0.85)$ -- estimates are within sampling
error. Persistence:
\[
  \widehat\alpha_1 + \widehat\beta_1 = 0.9029 < 1
  \quad\Rightarrow\quad
  \widehat\sigma^2 = \frac{0.0878}{1 - 0.9029}
                   = 0.9045,
\]
to be compared with the sample variance $0.9053$ -- excellent
match.

GARCH$(1, 1)$: $\log\mathcal L = -2009.31$, AIC $= 4026.62$.
ARCH$(5)$ has AIC $= 4028.54$. \textbf{GARCH$(1, 1)$ wins on AIC
with three parameters (instead of six)}: parsimony argument
confirmed.
\end{solution}

% =========================================================
\begin{exercise}[Standardised-residual diagnostics]\end{exercise}

\begin{solution}
On $\widehat z_t = \widehat\varepsilon_t / \widehat\sigma_t$:
\begin{center}
\begin{tabular}{rrrl}
\toprule
lag $q$ & LM stat (on $\widehat z_t^2$) & $p$-value & Decision \\
\midrule
$5$    & $4.67$  & $0.457$ & do not reject -- white \\
$10$   & $9.01$  & $0.531$ & do not reject -- white \\
$20$   & $15.61$ & $0.740$ & do not reject -- white \\
\bottomrule
\end{tabular}
\end{center}
The Engle LM no longer rejects at any lag: GARCH$(1, 1)$ has
absorbed the volatility clustering. The QQ-plot of
$\widehat z_t$ versus $\mathcal N(0, 1)$ is approximately linear,
consistent with the simulated Gaussian innovations.
\end{solution}

% =========================================================
\begin{exercise}[Forecast and Value-at-Risk]\end{exercise}

\begin{solution}
The closed-form mean-reversion identity
\[
  \widehat\sigma_{T+h}^2
  = \widehat\sigma^2
  + (\widehat\alpha_1 + \widehat\beta_1)^{h - 1}
        \bigl(\widehat\sigma_{T+1}^2 - \widehat\sigma^2\bigr)
\]
matches the \texttt{arch} package output to machine precision (max
discrepancy $6.7 \times 10^{-16}$). The figure
\texttt{lab6\_ex6\_forecast.png} shows the slow decay
toward $\widehat\sigma^2 = 0.9045$ at rate
$0.9029$.

\textbf{Value-at-Risk.} At $T$, the conditional sd is
$\widehat\sigma_{T+1} = 1.339$, larger than the unconditional sd
$0.951$: the series ends in a turbulent regime. Resulting
one-day VaR:
\begin{center}
\begin{tabular}{rrrr}
\toprule
$\alpha$ & $z_\alpha$ & VaR conditional & VaR unconditional \\
\midrule
$1\%$ & $-2.326$ & $3.07$ & $2.16$ \\
$5\%$ & $-1.645$ & $2.15$ & $1.52$ \\
\bottomrule
\end{tabular}
\end{center}
The conditional VaR is \emph{larger} -- correctly flagging higher
short-term risk that the static VaR misses.
\end{solution}

\end{document}
